% Data and likelihood notation.
Suppose we observe $N$ data points, $\xvec = (\x_1, \ldots \x_N)$.  We will
assume that the entries of $\xvec$ are independently and identically distributed
(IID) according to a distribution $\fdist$.  In a slight abuse of notation, we
will write $\xfdist$ for $\prod_{n=1}^N \fdist(\x_n)$.  Throughout, the
distribution $\fdist$ will be taken as fixed but unknown.  A Bayesian typically
models the distribution of $\xvec$ by an IID parametric model
given on some parameters $\theta \in \thetadom \subseteq \rdom{\thetadim}$:
%
$$
\p(\xvec \vert \theta) = \prod_{n=1}^N \p(\x_n \vert \theta) = 
    \exp\left(\sumn \ell(\x_n \vert \theta) \right)
\intertxt{where}
\ell(\x_n \vert \theta) := \log \p(\x_n \vert \theta).
$$
%
In the present work, we admit that the model may be misspecified, i.e., there may be no $\theta
\in \thetadom$ such that $\p(\x_n \vert \theta)$ is the same distribution as
$\fdist$.  

We will additionally posit a prior over the parameters
$\theta$, and assume this prior can be expressed as
a density $\prior(\cdot)$ with respect to the Lebesgue measure on $\rdom{\thetadim}$.  Then
we can compute the posterior expectation of a vector-valued quantity $g(\theta)
\in \rdom{\gdim}$ using Bayes' rule:
%
\begin{align}
%
\expect{\post}{\g(\theta)} =
    \frac{
    \int_{\thetadom} g(\theta)
        \exp\left(N \meann \ell(\x_n \vert \theta) \right) \prior(\theta)
                        d\theta}
         {\int_{\thetadom}
         \exp\left(N \meann \ell(\x_n \vert \theta) \right) \prior(\theta)
         d\theta}.
         \eqlabel{g_post}
%
\end{align}

The quantity $\expect{\post}{\g(\theta)}$ is a function of the data $\xvec$,
through its dependence on the random function $\theta \mapsto \meann \ell(\x_n
\vert \theta)$.  We will be primarily concerned with the problem of how to
estimate the variability of $\expect{\post}{\g(\theta)}$ under sampling from
$\xfdist$, when $N$ is large, using Markov Chain Monte Carlo (MCMC)
samples from the posterior $\post$, in the presence of possible
misspecification.
%
% Using an estimate of its variability, one can use classical frequentist tests to
% estimate how much the posterior mean might be expected to have varied were we to
% have drawn a different sample $\xvec$.
Specifically, let $\plim$ and $\dlim$ respectively denote convergence in
probability and distribution under sampling from $\xfdist$. We will
articulate conditions under which
% The mean and covariance (if they exist) of $\expect{\post}{\g(\theta)}$ under
% sampling from $\xfdist$ can be written respectively as
%
\begin{align}\eqlabel{gtrue}
%
\expect{\post}{\g(\theta)} \plim \gmeantrue < \infty
\quad\textrm{and}\quad
\sqrt{N}\left(\expect{\post}{\g(\theta)} - \gmeantrue \right)
\dlim
\normdist\left(\cdot \vert 0, \gcovtrue\right)
% \gcovtrue_{FS} := \cov{\xfdist}{\expect{\post}{\g(\theta)}}.
%
\end{align}
for some constants $\gmeantrue$ and $\gcovtrue$.  Under these conditions, our
contribution is to form a consistent estimator of $\gcovtrue$ that can be
computed from a single set of draws from $\post$, and so approximated from a
single run of MCMC.  We emphasize that our objective of establishing
\eqref{gtrue} is weaker than estimating the limiting variance of $\sqrt{N}
\expect{\post}{\g(\theta)}$, which may not even exist, but which would coincide
with $\gcovtrue$ under an additional assumption of uniform integrability (see,
e.g., Theorem 10.3.6 of \citet{dudley:2018:realanalysis}).




%%%%%%%%%%%%%%%%%%%%%%%%%%%%%%%%%%%%%
%%%%%%%%%%%%%%%%%%%%%%%%%%%%%%%%%%%%%

% Outline of the rest of the paper.


